% !TeX root = ../papers/nc-p3.tex
\section{異なる反対座標が強制する二十個の三重点を持つ八次曲面}
\label{sec:opposite-root-triple-octics}
\begin{theorem}\label{thm:opposite-root-triple-octics}
標準有限Frobenius窓の四辺が有限被約とし、ある位置の二十点を
$X\subset\mathbb P^3$ とする。その位置と反対位置の座標配置が標識を保って射影同値でなければ、
$H^0(\mathcal I_X^3(8))\ne0$ である。反対位置でも同じである。
実際その八次式は全二次式を成分として張る四行の対称二次式行列の行列式に取れ、
$X$ の各点で行列は階数1である。
\end{theorem}
\begin{proof}
定理\ref{thm:finite-square-chart}の正規化で、二つの二次Veronese像を同じ
$\mathbb P^9$ に置く。片側の $x\mapsto x^2$ を相手側の座標で読むと
対称四行行列 $M(x)$ となり、十個の二次式は基底である。$X$ で階数1となる。
一般階数4なら各点で単元を一つ消去して残りの三行の成分は極大イデアルに入るので
行列式は所望の八次式となる。一般階数1なら片方のVeronese三次元多様体が
他方に入って一致し、両表示は射影座標変更で同じになる。残る二階数を排除する。

一般階数2なら $\bigwedge^2M$、一般階数3なら $\operatorname{adj}M$ は
階数1の非零対称行列で、それぞれ次数4、6である。一意分解環でこの行列を
$h g g^{\mathsf t}$ と書く。実際分数体上の列直線の原始多項式ベクトル $g$ を取り、
各素元の付値で $h$ の分母がないことを確認すればよい。
$g$ の全成分の次数を $e$ とする。$e=0$ はどちらの場合も不可能である。
前者では二次元像が固定、後者では核直線が固定となり、行列の十成分の独立性に反する。

一般階数2では $1\le e\le2$、$\deg h=4-2e$ である。$X$ で全ての二次小行列式が消えるので、
$h g_i$ は全て $X$ 上で零となる。$X$ は被約であり、$4-e\le3$ 次の非零式が
$X$ を含むことは三次評価が同型であることに反する。

一般階数3では $1\le e\le3$、$\deg h=6-2e$ である。三次小行列式は $X$ で
二重以上に消える。$e=3$ なら $h$ は定数で全 $g_i$ が $X$ を含み矛盾する。
$e=2$ なら $g$ が零でない点では二次式 $h$ が二重に消える。
よって非零な $g_i\partial_j h$ を選ぶと全二十点を通る三次式になり矛盾する。

$e=1$ なら $g$ の共通零点は高々直線 $B$ である（原始性）。
$X\smallsetminus B$ は四次式 $h$ の特異点で、$X\cap B$ は高々三点である。
$h$ に四つの本質的変数があれば四つの独立な偏導三次式のうち、
非零な線型結合を残りの三点以下でも消せる。本質的変数が三つ以下なら $h$ を平面四次式の錐とみる。
その平面四次が非被約なら特異台は二次曲線に入る。実際重複二次曲線か、
重複直線と残余二次式の特異点を取ればよく、二本の重複直線も含まれる。
被約で特異点が五つ以下ならそれらを通る二次式がある。いずれもその錐の式に $g_i$ を掛ければ矛盾する。
残る被約平面四次の特異点数は高々6で、6の場合は四直線である。
正規化の種数公式 $\sum\delta_p=3+(\text{成分数})-1-\sum g_i$ による。
その四直線のうち三本ずつの積は四つの独立な三次式となり（各直線で制限する）、
全特異点で消える。再び $X\cap B$ に高々三条件を課せばよい。
以上で全ての可能性を尽くした。
\end{proof}

\begin{corollary}\label{cor:octic-free-finite-extension}
有限被約な第一辺を指定し、その最初の射影が八次の三重条件を満たす非零式を持たないとする。
四辺有限の標準正則延長は添字保持同型まで高々一つで、その有無は有限に決まる。
両射影でこの八次式がなければ全ての延長は両対称である。
さらに十五次元成分の閉包上では、非被約を含む全ての有限四辺の各位置で
$H^0(\mathcal I_X^3(8))\ne0$ が必要である。
\end{corollary}
\begin{proof}
第一の主張で四辺有限な延長があれば、第一辺が被約なため点列の同型で全て被約となる。
定理により反対二位置を同定でき、次の辺は第一辺の転置となる。
三次交叉の次元4を検査してその基底を最後の文字で束ねる。
命題\ref{prop:finite-flip-obstructions}により両端で対称であり、
標準有限Frobenius検査と四辺の有限性を通るものが唯一の候補である。
反対階数は10以下で、9以下なら系\ref{cor:all-defective-one-sym-finite}、
10なら定理\ref{thm:single-edge-instanton-extension}が必要十分判定を与える。
両射影に仮定すれば同じ議論で両交換が対称となる。
十五次元成分の一般点は有限被約で反対の二組がそれぞれ同定できない
（定理\ref{thm:veronese-pencil-quotient}、一般反対階数16）。
そこで各位置に上の非零八次式がある。有限辺の開集合上で各射影は
有限平坦長20であり、滑らかな相対三空間内の相対局所完全交叉である
（辺の評価階数3のところで三つの局所式を取る）。
従って三重化も平坦である。実際定義イデアルの次数1、2の随伴商は
余法束の対称冪であり、これらと構造層は底上平坦である。
八次評価の核の次元の上半連続性から閉包上の必要条件が従う。
\end{proof}

\paragraph{有限行列による検査。}
被約点での三重条件は二階偏導十本が全て $X$ を通ることと同値である
（標数零のEuler式で低階も従う）。各六次式を十五個の四次最大小行列式の二次式倍で割れば
$200\times165$ の行列になる。点を分離できる場合は直接その十本を評価してもよい。
\path{computations/verify_opposite_root_triple_octics.py} は一般の第一辺の整数標本と
対称完全交叉の既出の二十点例の両方で、良い素数で階数165を検証する。
整数標本の六次商次元20は期待余次元の最大小行列式分解による
（一般辺では有限被約であり、表示の階数が最大となる開条件を取る）。
よって一般辺でこの八次式は存在しない。また対称完全交叉のこの正則例は十五次元閉包の外にある。
その近傍では両側の八次条件の不在と被約性が開であるため、
そこを通る全ての正則な有限窓は両対称に限られる。
